% =============================================================================
% Limitations (mandatory unnumbered for *ACL since 2022; not counted in main pages)
% =============================================================================

\paragraph{Breadth-over-depth evaluation design.}
The main grid prioritises breadth across twelve
$(\text{harness}, \text{backbone}, \text{benchmark})$
combinations plus a frontier robustness check.  Depth is
uneven by design.  All three benchmarks are run with three
repeated evaluations per cell, and we report task-clustered
standard errors over the per-task paired deltas
(\S\ref{sec:main-results}).  On the high-baseline LLM-judged
benchmarks (GDPVal, QwenClawBench) per-task judge noise exceeds
the effect size, so the positive effect is established at the
pooled benchmark level rather than cell by cell.  SkillsBench
additionally admits paired per-task statistics (exact McNemar on
a representative run) and the logistic safety-facet regression
(Appendix~\ref{app:quality-detail}).  A larger
number of repetitions and seeds per cell would tighten the
per-cell estimates that the pooled analysis currently leans on.

\paragraph{Text-based skill scoring.}
The 3-facet quality framework and 19-flag taxonomy are evaluated
by an LLM judge against the \texttt{SKILL.md} text content rather
than by sandboxed execution.  This is a deliberate scale
trade-off.  Judge-time text scoring is feasible at the
$\sim$$821{,}000$-file crawl scale, whereas runtime sandboxing
would require per-skill compatible
execution environments not generally available for
community-contributed code.  Safety hard-gates and licence
filtering at Stage~5 catch the most severe content-level risks
at ingest time.  Complementary runtime verification at
retrieval-time or agent-execution-time is a principled extension
for future work.  A second consequence of the scale trade-off is
that the quality, safety, and classification labels driving the
release decision come from a single LLM judge and are not
validated against human gold annotations or cross-checked against
an alternative judge model.  Judge bias therefore propagates into
the released active set.  The facet-separation analysis is an internal-consistency
check, not a
construct-validity study, and human-annotation and
judge-robustness studies are left to future work.

\paragraph{Benchmark and harness coverage.}
The evaluation spans three publicly available benchmarks
(SkillsBench, GDPVal, QwenClawBench) chosen for
non-saturation on frontier LLMs and for complementary task-type
coverage (deterministic-verifier software tasks, real-world
knowledge work, and a
real-user-distribution agent benchmark), and two
\texttt{SKILL.md}-conformant agent harnesses (OpenClaw, Raven)
covering the most widely adopted skill-loading convention in
the public ecosystem.  Per-benchmark and per-harness
heterogeneity within this coverage is reported transparently
in the main results (\S\ref{sec:negative}) rather than
relegated to this section.
Extending the evaluation to additional task types
(e.g., long-horizon embodied tasks, non-English workflows) and
to harnesses with substantially different tool-use interfaces
or planning architectures would broaden external validity.  The
corpus is also English-dominant, both in its source repositories
and in our evaluation, which bounds claims about non-English
skills.  The
corpus-side findings (Stage~5 active--vs--filtered separation,
class-adaptive retrieval, and the suggestive safety-facet signal)
are expected to transfer modulo such adapter work, while per-cell
$\Delta$ magnitudes would re-anchor against each new harness's
own no-skill baseline.

\paragraph{Snapshot semantics.}
The released active set is a snapshot of the public
\texttt{SKILL.md} ecosystem at crawl time, while the community itself
continues to evolve (new repositories, licence updates, shifting
class distributions).  We do not study temporal churn or
re-curation cycles in this work, but the pipeline (per-stage
reject logs, deterministic Stage~5 cut, deduplication keyed on
content hash and embedding cosine) is constructed to support
periodic re-runs against future ecosystem snapshots.  The active
set released with this paper should be read as a 2026-Q2
snapshot rather than a one-shot artifact.
